% Gustavo Lazzari — DevOps / Site Reliability Engineer (SRE)
% Generated: 2026-02-26
% Template: Based on Rover Resume (https://github.com/subidit/rover-resume)

\documentclass[10pt]{article}

\usepackage[margin=0.6in, a4paper]{geometry}
\setcounter{secnumdepth}{0}
\usepackage{titlesec}
\titlespacing{\section}{0pt}{5pt}{2pt}
\titlespacing{\subsection}{0pt}{3pt}{1pt}
\titlespacing{\subsubsection}{0pt}{2pt}{1pt}
\titleformat{\section}{\large\bfseries\uppercase}{}{}{}[\titlerule]
\titleformat*{\subsubsection}{\normalsize\itshape}
\usepackage{enumitem}
\setlist[itemize]{noitemsep,topsep=1pt,left=0pt .. \parindent}
\setlist[description]{noitemsep,topsep=1pt,itemsep=1pt}
\pagestyle{empty}
\pdfgentounicode=1
\setlength{\parskip}{1pt}
\usepackage[utf8]{inputenc}
\usepackage[T1]{fontenc}
\usepackage{hyperref}
\hypersetup{hidelinks}
\urlstyle{same}

\begin{document}

\begin{center}
	{\LARGE\bfseries Gustavo Lazzari} \\
	\vspace{2pt}
	\href{mailto:gustavotlazzari@gmail.com}{gustavotlazzari@gmail.com} $|$
	+55 41 98773-2772 $|$
	\href{https://www.linkedin.com/in/glazzari}{linkedin.com/in/glazzari} $|$
	\href{https://github.com/GLazzari1428}{github.com/GLazzari1428} $|$
	Curitiba, PR, Brazil
\end{center}

\section{Education}
\subsection{Pontifical Catholic University of Paran\'{a} $|$ {\normalfont\itshape B.Sc. Computer Science} \hfill 2024 -- 2028}
\begin{itemize}
	\item Expected graduation: June 2028 $|$ Operating Systems, Network Security, Databases, Data Structures
\end{itemize}
\subsection{Swiss School of Curitiba $|$ {\normalfont\itshape High School \& IB Bilingual Diploma} \hfill 2017 -- 2020}

\section{Experience}
\subsection{Court of Justice of Paran\'{a} (TJPR) \hfill Curitiba, Brazil}
\subsubsection{DevOps / IT Intern \hfill Aug 2021 -- Dec 2022}
\begin{itemize}
	\item Designed and implemented CI/CD pipelines using GitHub Actions, Docker, and Docker Compose for automated testing, building, and deployment of web applications.
	\item Managed containerized application lifecycle across cloud platforms (AWS EC2, GCP Compute Engine), handling provisioning, deployment, and monitoring.
	\item Standardized development environments with Docker containers, reducing configuration drift between development and production deployments.
\end{itemize}

\section{Projects}
\subsection{Production-Grade Container Platform $|$ \normalfont\textit{20+ Services, Two-Node Cluster} \hfill Ongoing}
\begin{itemize}
	\item Architected and operate a two-node Proxmox VE cluster hosting 20+ containerized services (Docker/LXC), including media servers, NVR, private cloud, monitoring, and home automation --- treating personal infrastructure with production-level reliability standards.
	\item Engineered automated backup orchestration: dedicated secondary node manages all VM/container snapshots with offsite replication to Google Drive, ensuring disaster recovery readiness across the entire platform.
	\item Deployed high-availability DNS (PiHole across both nodes), eliminating single points of failure; configured ZFS storage with proactive disk health monitoring and hands-on failure recovery experience.
	\item Implemented GPU-accelerated transcoding pipeline (Quadro P400): PCI passthrough for Jellyfin VM and device mapping for Tdarr LXC, enabling distributed H.265 encoding across the cluster.
\end{itemize}

\subsection{Observability \& Secure Access Architecture $|$ \normalfont\textit{Monitoring, VPN, Zero Trust} \hfill Ongoing}
\begin{itemize}
	\item Deployed comprehensive monitoring stack: Prometheus for metrics collection, Grafana for visualization and alerting, and Scrutiny for predictive disk failure analysis via S.M.A.R.T. data.
	\item Engineered multi-layer access architecture: WireGuard VPN with Nginx reverse proxy for internal services, Cloudflare Zero Trust with TOTP for public-facing endpoints.
	\item Configured VLAN network segmentation (IoT/production) with OPNsense firewall rules; automated SSL/TLS certificate management via Certbot with Cloudflare DNS challenges.
	\item Manage identity infrastructure: self-hosted Vaultwarden for credential management and PocketID for centralized authentication.
\end{itemize}

\section{Technical Skills}
\begin{description}
	\item[Containers \& Virtualization] Docker, Docker Compose, LXC, Proxmox VE (multi-node), GPU passthrough, systemd
	\item[CI/CD \& Automation] GitHub Actions, Bash Scripting, YAML/IaC, Python, automated backup pipelines
	\item[Monitoring \& Observability] Prometheus, Grafana, Scrutiny (S.M.A.R.T.), log analysis
	\item[Networking \& Security] Nginx, WireGuard, Cloudflare (Tunnels/Zero Trust), OPNsense, PiHole, SSL/TLS, VLANs
	\item[Cloud \& Infrastructure] AWS (EC2), GCP (Compute Engine), Linux Admin (Arch/Debian), ZFS, RAID
	\item[Languages] Portuguese (Native), English (Fluent), German (Intermediate)
\end{description}

\section{Certifications}
Cambridge Advanced English (CAE -- B2) $|$ Deutsches Sprachdiplom (DSD -- B1/B2) $|$ Currently pursuing: AWS Certified Cloud Practitioner (CLF-C02)

\end{document}
